\section{One seed for the computational basis}\label{supp2:section:001}
\textbf{Manuscript item 2.4.}\quad \lean{ComputationalSeed.lean}\hfill 14 declarations

For $n\geq 0$, let
\[
 A_n=(\mathbb Z/2\mathbb Z)^{\{0,\ldots,n-1\}},\qquad
 V_n=\{v:A_n\to\C\}.
\]
For an address $a\in A_n$, define the standard coordinate vector
\[
 \ket a(x)=\begin{cases}1&x=a,\\0&x\ne a.\end{cases}
\]
The seed is $\psi_0=\ket0$. Define a translation operator by
\[
 (T_a v)(x)=v(x-a).
\]
At position $j$, the local Pauli $X$ is $X_j=T_{e_j}$, where $e_j$ has a single nonzero binary coordinate. The formal circuit uses the identity when $a_j=0$ and $X_j$ when $a_j=1$.

\begin{statement}\label{supp2:statement:001}
For every binary address $a$,
\[
 W_a=\prod_{j=0}^{n-1}X_j^{a_j}=T_a,
 \qquad W_a\psi_0=\ket a.
\]
The family $(W_a\psi_0)_{a\in A_n}$ is linearly independent, spans $V_n$, and has exactly $2^n$ elements. It is the standard computational basis.
\end{statement}

\begin{proof}\phantomsection\label{supp2:proof:001}
Composition gives $T_aT_b=T_{a+b}$, and direct evaluation gives $T_a\ket b=\ket{a+b}$. Every binary address satisfies $a+a=0$, so each translation is an involution. In particular, the local operator changes precisely the addressed bit:
\[
 X_j\ket a=\ket{a\text{ with coordinate }j\text{ replaced by }a_j+1}.
\]
The selected factor at $j$ equals translation by the address supported at $j$ with value $a_j$. Induction on a list proves that the product of translations is translation by their sum. Each of these addresses contributes one coordinate, and their sum is $a$. Thus $W_a=T_a$ and $W_a\psi_0=\ket a$.

Lean identifies the resulting family with \lean{Pi.basisFun} over $\C$. Its standard basis theorems supply linear independence and full span. The finite function-space cardinality gives $|A_n|=2^n$.
\phantomsection\label{supp2:proof:001:end}\end{proof}

\textbf{Formal scope.} The proof constructs the sequence of gates that produces each addressed state and then identifies the resulting family with the standard coordinate basis. It includes the empty circuit for $n=0$. Orthogonality belongs to the usual computational basis interpretation. The formal endpoint records the basis identification, together with separate explicit theorems for span and linear independence.





\section{Tensor and polyphase refinement}\label{supp2:section:002}
\textbf{Manuscript item 3.2.}\quad \lean{TensorRefinement.lean}\hfill 11 declarations

Let $F$ be an arbitrary finite set of factors and let $D_i$ be a finite coordinate index for factor $i$. For matrix families $A_i$, define the coordinate tensor product on the product index space by
\[
 \mathsf T(A)(x,y)=\prod_{i\in F} A_i(x_i,y_i).
\]
The proof first establishes the laws for multiplication, identity, adjoints, and finite sums. These laws apply to any finite number of factors with the stated finite coordinate sets.

Suppose $N>0$ and each factor has a complete orthogonal self-adjoint projector family $(P^{(i)}_j)_{j\in\mathbb Z/N\mathbb Z}$. For a fine address $a:F\to\mathbb Z/N\mathbb Z$, set
\[
 Q_a=\mathsf T\bigl(i\mapsto P^{(i)}_{a_i}\bigr),
 \qquad
 C_r=\sum_{\sum_i a_i=r}Q_a.
\]

\begin{statement}\label{supp2:statement:002}
The fine and coarse families are complete orthogonal self-adjoint projector systems. In particular,
\[
 C_rC_s=\delta_{rs}C_r,\qquad
 \sum_r C_r=I,\qquad C_r^*=C_r.
\]
Each coarse projector is exactly the displayed sum of its fine descendants.
\end{statement}

\begin{proof}\phantomsection\label{supp2:proof:002}
Expanding the finite product of sums gives
\[
 \mathsf T(A)\mathsf T(B)=\mathsf T(i\mapsto A_iB_i).
\]
Consequently, if two fine addresses differ in any coordinate, the corresponding product contains a zero factor. If the addresses agree, the product is the same fine projector. The adjoint passes through the finite coordinate product, so $Q_a^*=Q_a$. Summing over all address functions and using the same expansion yields
\[
 \sum_a Q_a=\mathsf T\left(i\mapsto\sum_jP^{(i)}_j\right)=I.
\]
A separate algebraic regrouping theorem applies to any finite orthogonal family and any address map into a finite set. In the double sum for $C_rC_s$, orthogonality removes all terms with different fine addresses. A surviving address can lie in both classes only when $r=s$. Completeness follows by summing each fine projector exactly once. Adjoint preservation follows term by term. Specializing the address map to total residue proves the claim.
\phantomsection\label{supp2:proof:002:end}\end{proof}

\textbf{Formal scope.} Residue grouping is proved for every positive modulus and arbitrary finite factor dimensions. It assumes the supplied local projector laws. Constructing those local projectors from cyclic characters, as in item 3.1, requires the separate proof in the next supplement.





\section{Differentiated moving coordinates}\label{supp2:section:003}
\textbf{Manuscript item 6.1.}\quad \lean{MovingFrames.lean}\hfill 9 declarations

Fix a real time $t$ and a finite complex coordinate space. Let $U(s)$ be a unitary matrix curve, differentiable entrywise at $t$, and let $\psi(s)$ be differentiable componentwise at $t$. Write $\dot U$ and $\dot\psi$ for those derivatives. Suppose
\[
 i\dot\psi=H\psi(t).
\]
Only the value of the Hamiltonian at the chosen time is needed. No derivative of $H$ is assumed.

\begin{statement}\label{supp2:statement:003}
The explicitly defined coordinate curve
\[
 y(s)=U(s)^*\psi(s)
\]
is differentiable at $t$, with
\[
 \dot y=\dot U^*\psi(t)+U(t)^*\dot\psi.
\]
It satisfies
\[
 \boxed{\quad i\dot y=H_{\mathrm{eff}}y(t),\qquad
 H_{\mathrm{eff}}=U(t)^*HU(t)-iU(t)^*\dot U.\quad}
\]
The connection term $-iU(t)^*\dot U$ is Hermitian. If $H^*=H$, then $H_{\mathrm{eff}}^*=H_{\mathrm{eff}}$.
\end{statement}

\begin{proof}\phantomsection\label{supp2:proof:003}
Applying the product rule entry by entry proves the differentiation formulas for matrix products and for matrices acting on vectors. Since the parameter is real, conjugation commutes with differentiation, giving the stated derivative of $U^*$ and hence of $y$.

Differentiate $U^*U=I$ using those rules and uniqueness of derivatives:
\[
 \dot U^*U+U^*\dot U=0.
\]
Thus $(U^*\dot U)^*=-U^*\dot U$, so multiplication by $-i$ produces a Hermitian matrix.

The identity $UU^*=I$ gives $\psi=Uy$ for the constructed coordinate curve. Differentiating and inserting the original equation yields
\[
 i\dot U y+iU\dot y=HUy.
\]
Multiplying by $U^*$ and using $U^*U=I$ gives the transformed equation, with the negative sign on the connection term. Finally, conjugation by $U$ preserves Hermiticity of $H$, and the sum of the two Hermitian terms is Hermitian.
\phantomsection\label{supp2:proof:003:end}\end{proof}

\textbf{Formal scope.} The main certificate constructs $y$ and its derivative rather than assuming a separate derivative witness for the desired answer. Both unitary identities are explicit. The proofs are pointwise and apply at any time where the stated hypotheses hold. They do not assert existence of a global solution to an evolution equation.





\section{The canonical packet connection}\label{supp2:section:004}
\textbf{Manuscript item 6.3.}\quad \lean{CanonicalConnection.lean}\hfill 9 declarations

Let $P_a(s)$ be a finite family in a normed real algebra with a continuous compatible adjoint. Assume actual derivatives $D_a=\dot P_a(t)$ exist and, for all parameter values under consideration,
\[
 P_aP_b=\delta_{ab}P_a,\qquad \sum_aP_a=I,\qquad P_a^*=P_a.
\]
In the supplied Lean theorem, the curves are defined on $\R$ and the projector laws hold at every real parameter. The algebraic formulation applies to matrices and bounded operators whenever they satisfy the stated typeclass hypotheses.

\begin{statement}\label{supp2:statement:004}
Define $C=\sum_aD_aP_a(t)$. Then
\[
 \boxed{\quad C=\frac12\sum_a\comm{D_a}{P_a(t)},\qquad C^*=-C,\quad}
\]
and, for every $b$,
\[
 \boxed{\comm{C}{P_b(t)}=D_b.}
\]
\end{statement}

\begin{proof}\phantomsection\label{supp2:proof:004}
The derivative product rule and uniqueness of derivatives give
\[
 D_aP_b+P_aD_b=\delta_{ab}D_a.
\]
Differentiating completeness and adjoint invariance gives
\[
 \sum_aD_a=0,\qquad D_a^*=D_a.
\]
These are proved derivative consequences, not additional axioms supplied in place of differentiation.

On the diagonal, $D_aP_a+P_aD_a=D_a$. Summing shows
\[
 \sum_aP_aD_a=-\sum_aD_aP_a=-C.
\]
This proves the formula involving half the sum of commutators, together with $C^*=-C$.

Orthogonality directly yields $CP_b=D_bP_b$. For the left product, multiply the differentiated identity for $(b,a)$ on the right by $P_a$ and sum over $a$. Completeness reduces the result to
\[
 D_b+P_bC=D_bP_b.
\]
Subtracting this equation from the identity for multiplication on the right gives $CP_b-P_bC=D_b$.
\phantomsection\label{supp2:proof:004:end}\end{proof}

\textbf{Formal scope.} The proof derives the differentiated identities before using their algebraic consequences, so the role of the derivative hypotheses remains explicit. It establishes the canonical connection identity for one parameter. The range bundle over a parameter manifold, the exterior differential calculus, and the curvature theorem of item 6.4 are developed in the later supplements.





\section{Finite Kraus channels and full amplification}\label{supp2:section:005}
\textbf{Manuscript item 7.2.}\quad \lean{KrausChannels.lean}\hfill 7 declarations

For a finite family of complex matrices $K_a$, define
\[
 \Phi(\rho)=\sum_a K_a\rho K_a^*.
\]
The formalization constructs this as a complex-linear map and proves the following properties directly in the finite matrix representation.

\begin{statement}\label{supp2:statement:005}
The map $\Phi$ is completely positive. If
\[
 \sum_aK_a^*K_a=I,
\]
then $\tr\Phi(\rho)=\tr\rho$ for every matrix $\rho$. In particular, positive semidefinite matrices of trace one are mapped to positive semidefinite matrices of trace one.
\end{statement}

\begin{proof}\phantomsection\label{supp2:proof:005}
Matrix multiplication distributes across addition and scalar multiplication, which gives complex linearity. If $\rho\succeq0$, each term $K_a\rho K_a^*$ is positive semidefinite, hence so is their finite sum.

To prove complete positivity rather than positivity alone, let $J$ be an arbitrary finite auxiliary index set. For a matrix $X$ indexed by $J\times n$, define its blockwise amplification by
\[
 [\Phi^{(J)}(X)]_{(i,u),(j,v)}
 =\Phi\bigl(X_{(i,\cdot),(j,\cdot)}\bigr)_{u,v}.
\]
The explicit index calculation proves the identity
\[
 \boxed{\quad
 \Phi^{(J)}(X)=\sum_a (I_J\otimes K_a)X(I_J\otimes K_a)^*.
 \quad}
\]
The displayed identity identifies the defined amplification with tensoring by the identity, which is the operation needed for complete positivity. Every summand on the right is positive when $X$ is positive, and this holds for every finite $J$.

Finally, cyclicity and linearity of the trace give
\begin{align*}
 \tr\Phi(\rho)
 &=\sum_a\tr(K_a\rho K_a^*)\\
 &=\tr\left(\left(\sum_aK_a^*K_a\right)\rho\right)
 =\tr\rho.
\end{align*}
Positivity together with trace preservation gives the stated conclusion for density matrices.
\phantomsection\label{supp2:proof:005:end}\end{proof}

\textbf{Formal scope.} Complete positivity is proved for every finite auxiliary dimension. Normalization is required only for trace preservation and the conclusion about density matrices. The converse representation theorem for arbitrary completely positive maps lies outside this argument.





